\documentclass[aspectratio=169]{beamer}
\usepackage[utf8]{inputenc}
\usepackage[T1]{fontenc}
\usepackage{lmodern}
\usepackage{xcolor}
\usepackage[portuguese]{babel}

\definecolor{BgCream}{HTML}{FAF8F5}
\definecolor{TextEspresso}{HTML}{4A4238}
\definecolor{NudeTaupe}{HTML}{BFAFA6}
\definecolor{SoftBlush}{HTML}{D9C5B2}

\setbeamertemplate{navigation symbols}{}
\setbeamercolor{background canvas}{bg=BgCream}
\setbeamercolor{normal text}{fg=TextEspresso}
\setbeamercolor{structure}{fg=TextEspresso}

\setbeamertemplate{itemize item}{\color{NudeTaupe}\textendash}
\setbeamertemplate{itemize subitem}{\color{NudeTaupe}\(\circ\)}

\setbeamerfont{title}{size=\huge, series=\bfseries}
\setbeamerfont{frametitle}{size=\Large, series=\bfseries}

\setbeamertemplate{title page}{
  \vspace*{1.5cm}
  \begin{center}
    \usebeamerfont{title}\textcolor{TextEspresso}{\inserttitle}\par
    \vspace{0.35cm}
    {\large\textcolor{NudeTaupe}{\insertsubtitle}}\par
    \vspace{0.6cm}
    {\color{NudeTaupe}\rule{2.5cm}{0.5pt}}\par
    \vspace{0.8cm}
    \usebeamerfont{author}\textcolor{TextEspresso}{\insertauthor}\par
    \vspace{0.3cm}
    \usebeamerfont{institute}\small\textcolor{NudeTaupe}{\insertinstitute \quad|\quad \insertdate}
  \end{center}
}

\setbeamertemplate{frametitle}{
  \vspace*{0.8cm}
  \usebeamerfont{frametitle}\insertframetitle\par
  \vspace*{-0.1cm}
  {\color{NudeTaupe}\rule{1.5cm}{0.5pt}}
  \vspace*{0.2cm}
}

\newenvironment{ideiachave}{
  \vspace{0.5cm}
  \begin{center}
  \color{TextEspresso}
  \itshape\Large
}{
  \end{center}
  \vspace{0.5cm}
}

\usepackage{csquotes}

\title{Contentores e Isolamento de Processos}
\subtitle{Infraestruturas de Computação na Nuvem, Aula 04}
\author{Catarina Silva}
\institute{ESTGA, Universidade de Aveiro}
\date{2026}

\begin{document}

\begin{frame}[plain]
  \titlepage
\end{frame}

\begin{frame}
\frametitle{Agenda \textcolor{NudeTaupe}{\small(1/2)}}
\begin{itemize}
    \setlength\itemsep{0.4cm}
    \item Contentores vs máquinas virtuais.
    \item Primitivas Linux de isolamento: namespaces e cgroups.
    \item Union filesystems e imagens em camadas.
\end{itemize}
\end{frame}

\begin{frame}
\frametitle{Agenda \textcolor{NudeTaupe}{\small(2/2)}}
\begin{itemize}
    \setlength\itemsep{0.4cm}
    \item Arquitetura Docker: daemon, CLI, containerd, runc.
    \item Registries de imagens e versionamento por tags.
    \item Redes e volumes em Docker.
    \item Boas práticas e segurança de contentores.
    \item Contentores de aplicação vs contentores de sistema.
    \item OCI: portabilidade entre runtimes.
\end{itemize}
\end{frame}

\begin{frame}
\frametitle{Contentores vs Máquinas Virtuais}
\begin{itemize}
    \setlength\itemsep{0.4cm}
    \item Contentores partilham o kernel do sistema operativo do host, ao contrário das VMs, que virtualizam hardware e correm um kernel próprio.
    \item Overhead muito menor: sem hipervisor completo nem sistema operativo convidado duplicado.
    \item Arranque quase instantâneo, tipicamente em segundos (ou menos), em vez dos minutos típicos de uma VM.
    \item Contrapartida: isolamento menos forte do que numa VM, já que dependem inteiramente das primitivas de isolamento do kernel do host.
\end{itemize}
\end{frame}

\begin{frame}
\frametitle{Namespaces Linux \textcolor{NudeTaupe}{\small(1/2)}}
\begin{itemize}
    \setlength\itemsep{0.4cm}
    \item Namespaces dão a cada processo uma visão isolada de um determinado recurso do sistema, criando a ilusão de máquina dedicada.
    \item \textbf{PID}: árvore de processos própria, isolada da do host.
    \item \textbf{NET}: pilha de rede própria (interfaces, rotas, portas).
    \item \textbf{MNT}: sistema de ficheiros e pontos de montagem próprios.
\end{itemize}
\end{frame}

\begin{frame}
\frametitle{Namespaces Linux \textcolor{NudeTaupe}{\small(2/2)}}
\begin{itemize}
    \setlength\itemsep{0.4cm}
    \item \textbf{UTS}: hostname e domainname próprios.
    \item \textbf{IPC}: mecanismos de comunicação entre processos isolados (semáforos, filas de mensagens).
    \item \textbf{USER}: mapeamento de identificadores de utilizador e grupo, permitindo root dentro do contentor sem privilégios no host.
    \item Um contentor Docker é, na prática, um processo Linux normal ao qual foram atribuídos vários destes namespaces.
\end{itemize}
\end{frame}

\begin{frame}
\frametitle{Cgroups: Limitação de Recursos}
\begin{itemize}
    \setlength\itemsep{0.4cm}
    \item Enquanto os namespaces isolam o que um processo \emph{vê}, os cgroups (control groups) limitam o que um processo pode \emph{consumir}.
    \item Limitação de CPU: quota e período de utilização atribuídos a um grupo de processos.
    \item Limitação de memória: teto máximo, com possibilidade de terminar o processo (OOM kill) se for excedido.
    \item Limitação de I/O: controlo de largura de banda de leitura/escrita em disco por grupo.
\end{itemize}
\end{frame}

\begin{frame}
\frametitle{Cgroups: Versões e Funcionamento}
\begin{itemize}
    \setlength\itemsep{0.4cm}
    \item Os cgroups são expostos pelo kernel através de um sistema de ficheiros virtual, tipicamente em \texttt{/sys/fs/cgroup}.
    \item Escrever valores nesses ficheiros define os limites aplicados a um grupo de processos.
    \item \textbf{cgroups v1}: hierarquias separadas e independentes para cada tipo de recurso (CPU, memória, I/O).
    \item \textbf{cgroups v2}: hierarquia única e unificada, com um modelo de delegação mais coerente; é o modelo usado por omissão nas distribuições Linux recentes.
    \item Sem cgroups, um único contentor descontrolado poderia consumir toda a memória ou CPU do host.
\end{itemize}
\end{frame}

\begin{frame}
\frametitle{Isolamento Não é Segurança Absoluta}
\begin{itemize}
    \setlength\itemsep{0.4cm}
    \item Todos os contentores de um host partilham o mesmo kernel: uma vulnerabilidade no kernel é uma superfície de ataque comum.
    \item \textbf{Capabilities}: o Linux divide os privilégios de root em capacidades individuais (ex.\ \texttt{CAP\_NET\_ADMIN}); um contentor deve receber apenas as estritamente necessárias.
    \item \textbf{Seccomp}: filtra as chamadas de sistema (\textit{syscalls}) que um processo pode invocar, reduzindo a superfície exposta do kernel.
    \item \textbf{AppArmor} e \textbf{SELinux}: mecanismos de controlo de acesso obrigatório, que restringem o que cada processo pode aceder.
    \item Quando é exigido isolamento forte (ex.\ entre clientes diferentes), uma VM continua a ser mais segura do que um contentor.
\end{itemize}
\end{frame}

\begin{frame}
\frametitle{Union Filesystems e Imagens em Camadas}
\begin{itemize}
    \setlength\itemsep{0.4cm}
    \item OverlayFS combina múltiplos diretórios (camadas) numa única vista de sistema de ficheiros.
    \item Uma imagem Docker é composta por várias camadas só de leitura, empilhadas umas sobre as outras.
    \item Ao correr um contentor, é adicionada uma camada de escrita (writable layer) no topo, específica dessa instância.
    \item Copy-on-write: alterações a ficheiros de camadas inferiores são copiadas para a camada de escrita antes de modificadas, preservando a imagem original.
\end{itemize}
\end{frame}

\begin{frame}
\frametitle{Arquitetura Docker \textcolor{NudeTaupe}{\small(1/2)}}
\begin{itemize}
    \setlength\itemsep{0.4cm}
    \item \textbf{Docker CLI}: interface de linha de comandos usada pelo utilizador.
    \item \textbf{dockerd}: daemon que recebe os pedidos da CLI (ou da API) e gere imagens, redes e volumes.
    \item \textbf{containerd}: runtime de alto nível, responsável pelo ciclo de vida dos contentores (criação, execução, paragem).
    \item \textbf{runc}: runtime de baixo nível que efetivamente cria o contentor, aplicando namespaces e cgroups.
\end{itemize}
\end{frame}

\begin{frame}
\frametitle{Arquitetura Docker \textcolor{NudeTaupe}{\small(2/2)}}
\begin{itemize}
    \setlength\itemsep{0.4cm}
    \item Um Dockerfile descreve, em passos declarativos, como construir uma imagem (imagem base, dependências, código, comando de arranque).
    \item Cada instrução do Dockerfile gera, em regra, uma nova camada da imagem.
    \item O comando \texttt{docker build} executa estes passos e produz a imagem final.
    \item A imagem resultante pode depois ser instanciada em um ou mais contentores com \texttt{docker run}.
\end{itemize}
\end{frame}

\begin{frame}
\frametitle{Registries de Imagens}
\begin{itemize}
    \setlength\itemsep{0.4cm}
    \item Um registry é um repositório de imagens Docker, de onde estas podem ser descarregadas (\texttt{pull}) ou para onde podem ser enviadas (\texttt{push}).
    \item Docker Hub é o registry público de referência, com imagens oficiais e da comunidade.
    \item Organizações usam também registries privados, para imagens internas ou proprietárias.
    \item Versionamento por tags: identificam versões específicas de uma imagem (ex.: \texttt{nginx:1.25}, \texttt{minha-app:latest}).
\end{itemize}
\end{frame}

\begin{frame}
\frametitle{Redes em Docker}
\begin{itemize}
    \setlength\itemsep{0.4cm}
    \item \textbf{bridge}: rede virtual isolada no host, usada por omissão; os contentores comunicam entre si e saem para o exterior por NAT.
    \item Numa rede bridge \emph{personalizada}, o Docker fornece resolução de nomes: um contentor alcança outro pelo seu nome.
    \item \textbf{host}: o contentor usa diretamente a pilha de rede do host, sem isolamento de rede (sem namespace NET próprio).
    \item \textbf{none}: sem qualquer interface de rede, para cargas de trabalho totalmente isoladas.
    \item Publicação de portas (\texttt{-p 8080:80}) mapeia uma porta do host para uma porta do contentor.
\end{itemize}
\end{frame}

\begin{frame}
\frametitle{Persistência: Volumes e Bind Mounts}
\begin{itemize}
    \setlength\itemsep{0.4cm}
    \item A camada de escrita de um contentor desaparece quando este é removido -- por omissão, os dados não persistem.
    \item \textbf{Volumes}: geridos pelo Docker, guardados numa área própria do host; a forma recomendada de persistir dados.
    \item \textbf{Bind mounts}: mapeiam um diretório concreto do host para dentro do contentor; úteis em desenvolvimento, mas dependentes da estrutura do host.
    \item \textbf{tmpfs}: armazenamento apenas em memória, que nunca chega ao disco -- adequado a dados sensíveis ou temporários.
    \item Separar dados do ciclo de vida do contentor é essencial para tratar contentores como recursos descartáveis.
\end{itemize}
\end{frame}

\begin{frame}
\frametitle{Boas Práticas na Construção de Imagens}
\begin{itemize}
    \setlength\itemsep{0.4cm}
    \item Usar imagens base minimalistas (ex.\ Alpine, ou variantes \textit{slim}) para reduzir tamanho e superfície de ataque.
    \item \textbf{Multi-stage builds}: compilar numa fase intermédia e copiar apenas o resultado final para a imagem entregue.
    \item Ordenar as instruções do Dockerfile do menos ao mais volátil, para aproveitar a cache de camadas (ex.\ instalar dependências antes de copiar o código).
    \item \texttt{.dockerignore}: evita enviar ficheiros desnecessários para o contexto de build.
    \item Não incluir segredos (passwords, chaves) na imagem: ficam registados nas camadas e são recuperáveis.
    \item Correr a aplicação com um utilizador sem privilégios, em vez de root.
\end{itemize}
\end{frame}

\begin{frame}
\frametitle{Contentores de Aplicação vs de Sistema}
\begin{itemize}
    \setlength\itemsep{0.4cm}
    \item \textbf{Contentor de aplicação} (Docker): corre tipicamente um único processo/serviço; é efémero e descartável.
    \item Não tem um sistema de init completo; o ciclo de vida do contentor é o ciclo de vida desse processo.
    \item \textbf{Contentor de sistema} (LXC, visto na Aula 03): comporta-se como uma máquina Linux completa, com \texttt{systemd} e múltiplos serviços.
    \item Ambos usam as mesmas primitivas do kernel (namespaces e cgroups) -- a diferença está na forma de utilização, não na tecnologia de base.
    \item Nos laboratórios desta UC usamos os dois: LXC no Proxmox para máquinas de sistema, Docker para empacotar aplicações.
\end{itemize}
\end{frame}

\begin{frame}
\frametitle{OCI: Open Container Initiative}
\begin{itemize}
    \setlength\itemsep{0.4cm}
    \item Iniciativa que define especificações abertas e neutras de fornecedor para formatos de contentores.
    \item \textbf{image-spec}: define o formato de uma imagem de contentor, independente da ferramenta que a construiu.
    \item \textbf{runtime-spec}: define como um runtime deve criar e executar um contentor a partir dessa imagem.
    \item Garante portabilidade: a mesma imagem pode correr em diferentes runtimes compatíveis com OCI, como runc, crun ou gVisor.
\end{itemize}
\end{frame}

\begin{frame}
\frametitle{Síntese da Aula \textcolor{NudeTaupe}{\small(1/2)}}
\begin{itemize}
    \setlength\itemsep{0.4cm}
    \item Contentores partilham o kernel do host, oferecendo menor overhead e arranque quase instantâneo face a VMs.
    \item Namespaces isolam o que um processo vê (PID, NET, MNT, UTS, IPC, USER); cgroups limitam o que consome (CPU, memória, I/O).
    \item OverlayFS e imagens em camadas permitem reutilização eficiente de camadas comuns entre imagens.
\end{itemize}
\end{frame}

\begin{frame}
\frametitle{Síntese da Aula \textcolor{NudeTaupe}{\small(2/2)}}
\begin{itemize}
    \setlength\itemsep{0.4cm}
    \item Docker articula dockerd, containerd e runc para transformar um Dockerfile numa imagem e, depois, num contentor em execução.
    \item Registries (como o Docker Hub) distribuem imagens versionadas por tags; redes bridge e volumes tratam de comunicação e persistência.
    \item O isolamento de contentores não é absoluto: capabilities, seccomp e AppArmor/SELinux reforçam-no, mas uma VM continua a isolar melhor.
    \item OCI garante portabilidade das imagens entre diferentes runtimes de contentores.
\end{itemize}
\end{frame}

\begin{frame}[plain]
  \vspace*{3cm}
  \begin{center}
    \Huge\bfseries\textcolor{TextEspresso}{Obrigada.}\par
    \vspace{0.5cm}
    {\color{NudeTaupe}\rule{2cm}{0.5pt}}\par
    \vspace{0.5cm}
    \Large\textcolor{TextEspresso}{\itshape Questões e comentários}
  \end{center}
\end{frame}

\end{document}
